\section{Introduction}
A finite stochastic register has a continuum of probability distributions. Its number of labels is therefore not determined by the cardinality of the numerical orbit it realizes. Allowing an external clock creates a further difficulty: the command rows may change at every position, and the entire realization may be redesigned at each horizon. The central question is whether these freedoms change the least width needed to preserve a fixed numerical experiment with a prescribed uniform error.

Our first theorem identifies the limiting clocked width with the stationary minimum, without a loss of labels or accuracy. The synchronization hypothesis is eventual approximation of the physical identity at every sufficiently large length. Under it, arbitrarily many cuts of width at most $k$ force a stationary realization on $k$ labels, even when the intervening registers are unrestricted. The proof selects approximately homogeneous composite transition tables by finite Ramsey theory. It applies to irreversible topological monoids as well as groups. It does not assume that a physical return acts trivially on the hidden register.

For compact groups a second reduction makes the stationary realization physically equivariant. Transition randomness can be removed, and hidden permutations above the physical identity can be averaged out, without increasing width or uniform error. The stationary minimum is therefore intrinsic to a finite continuous group action with randomized initialization. Combining the two reductions gives an equality-inclusive finite-quotient criterion for clocked boundedness, including positive critical radii on groups with infinitely many components. A moving-frame construction gives the exact limiting minimum in each return-length phase, again without additional labels.

\subsection{The interface and its two resources}
Unless a section explicitly specifies a different physical space, $H$ is a compact Hausdorff group and $\mathcal A$ is a finite nonempty alphabet whose letters $u_a\in H$ generate a dense subgroup. The execution convention is
\[
                  u_{a_1\cdots a_n}=u_{a_n}\cdots u_{a_1}.
\]
There are finite nonempty seed and query sets $\mathcal S,\mathcal J$ and continuous maps $F_{s,j}:H\to Y_j$. In the structural results $Y_j$ is a nonempty norm-compact convex subset of a real Banach space. The localization and algebraic sections specify finite-dimensional output spaces. For a categorical law, $Y_j=\Delta_{M_j}$ and the norm on differences is $\frac12\|\cdot\|_1$, namely total variation. A prescribed joint law is one query; incompatible measurements are not assigned an artificial joint law.

\begin{definition}[Stationary realization]\label{def:stationary55}
A realization on $k$ labels consists of initialization rows $\alpha_s\in\Delta_k$, row-stochastic matrices $T_a$ of order $k$, and decoder values $D_j(i)\in Y_j$, all independent of time and word length. Its error is
\[
       \sup_{s,j,w\in\mathcal A^*}
        \|\alpha_sT_{a_1}\cdots T_{a_n}D_j-F_{s,j}(u_w)\|_j.
\]
The empty word is included. Let $M_\varepsilon$ be the least $k$ at error at most $\varepsilon$, or $\infty$ when none exists. A permutation realization has permutation command matrices; its initialization need not be deterministic.
\end{definition}
The clocked resource $W_{N,\varepsilon}$ instead minimizes the maximum size of the registers $S_0,\ldots,S_N$ at a fixed horizon, with arbitrary row-stochastic transitions depending on time and horizon, and a terminal query decoder. Definition~\ref{def:machine54} states the model formally. Every advertised persistent label counts, including unused padding. Private information retained for future commands is part of the current label. The original width model does not charge the external clock, real table descriptions, their construction, exact arithmetic, or exact atomic sampling. Section~\ref{sec:complexity56} gives separate finite-description and fair-random-bit bounds; it does not reinterpret those resources as free computational memory. A numerical decoder vector is not an additional sampled-answer register.

For an open normal subgroup $U\le H$ put
\begin{equation}\label{eq:quotientradius55}
 r(U)=\max_{s,j,h\in H}\rad_{Y_j}(F_{s,j}(hU)),
 \qquad\rad_Y(A)=\min_{y\in Y}\sup_{x\in A}\|x-y\|.
\end{equation}
The coset images are compact by continuity. The radius depends continuously on the translate, and compactness gives the displayed extrema. The restriction of the center to the legal output set is essential. The component radius $R_0$ uses $U=H^\circ$, whether or not this subgroup is open; it is the infimum of the finite-quotient radii. At the critical error, the relevant issue is attainment of that infimum, not a limit from larger errors.

\subsection{The principal conclusions}
Theorem~\ref{thm:stationarization56} proves $W_{N,\varepsilon}\to M_\varepsilon$ under eventual approximate returns. Finite limits are eventually equal, and every $k<M_\varepsilon$ has a uniform narrow-cut occupation bound. Theorems~\ref{thm:equivariant56} and \ref{thm:completeboundary56} then give the intrinsic group-action minimum and its finite-quotient existence criterion. Theorem~\ref{thm:phaseminimum56} treats each length phase by a time-dependent change of physical frame. Table~\ref{tab:map56} separates these quantifiers and hypotheses.

The quantitative consequences have additional hypotheses. For a specified Diophantine planar alphabet the full-subcritical matching polynomial law is retained. A fixed $2$-adic interface has linear exact width and bounded width at every positive error. For a spherical interface in dimension $d$ with a spectral-gap alphabet, Theorem~\ref{thm:spherical56} gives the order $N^{(d-1)/2}$ up to a logarithmic factor in the lower bound. A rational alphabet generating $\mathrm{SO}(3)$ supplies a genuinely noncommuting example. Its spectral gap is imported from Bourgain--Gamburd, not inferred from a finite numerical test. Theorem~\ref{thm:ETR56} gives an $\exists\R$-complete finite input problem, while Theorem~\ref{thm:bits56} prices dyadic descriptions and fair random bits.

The compact nonnegative-matrix group structure used in purification is classical, and is separated from the external numerical error constraints throughout Section~\ref{sec:comparison54}. Neutral-letter collapse for advice automata is also a relevant predecessor to stationarization. The additional assertions proved here concern the same stochastic label bound and the same closed numerical error balls, arbitrary intervening widths, and approximate rather than exact returns. Neither classical theorem is relabeled as a new result.

\clearpage
\begin{table}[p]
\centering\small
\caption{Resources, hypotheses, and conclusions. All widths count persistent labels; all errors are worst-word numerical errors.}\label{tab:map56}
\begin{tabular}{p{.20\textwidth}p{.35\textwidth}p{.35\textwidth}}
\toprule
Object & Hypotheses & Conclusion\\
\midrule
$M_\varepsilon$ & Compact $H$; finite dense alphabet; compact convex outputs; no returns needed & Least finite continuous $H$-action with randomized initialization (Theorem~\ref{thm:equivariant56})\\[5pt]
$r(U)$ & $U$ open normal; arbitrary number of components & $M_\varepsilon<\infty$ iff some $r(U)\le\varepsilon$, equality included (Theorem~\ref{thm:finitequotient55})\\[5pt]
$W_{N,\varepsilon}$ & Eventual approximate returns; even an irreversible topological monoid & $W_{N,\varepsilon}\to M_\varepsilon$; finite limits eventually equal; all-profile occupation (Theorem~\ref{thm:stationarization56})\\[5pt]
$R_0$ & Compact group and eventual approximate returns & $R_0=\inf_Ur(U)$; critical boundedness iff the infimum is attained (Theorem~\ref{thm:completeboundary56})\\[5pt]
$M_{r,\varepsilon}$ & Block length $p$, phase subgroup $J$; block approximate returns & $W_{N,\varepsilon}\to M_{r,\varepsilon}$ for $N\equiv r$; no extra phase labels (Theorem~\ref{thm:phaseminimum56})\\[5pt]
$\varepsilon_r$ & Finitely many components; nonempty exact-return language & Component radius in the reachable phase; earlier localized occupation proof (Theorem~\ref{thm:phases55})\\[5pt]
Exact boundary & Cofinite exact returns; finite-dimensional outputs & Finite translate rank yields a finite observable quotient (Theorem~\ref{thm:zero55})\\[5pt]
Quantitative width & Separate Diophantine, profinite, or spherical spectral-gap hypotheses & Matching planar and dyadic orders; spherical order up to a logarithm (Theorems~\ref{thm:matched54}, \ref{thm:adiclinear55}, \ref{thm:spherical56})\\[5pt]
Finite encodings & Explicit finite group and rational categorical tables & $\exists\R$-complete stationary decision; dyadic precision and random-bit budgets (Theorems~\ref{thm:ETR56}, \ref{thm:bits56})\\
\bottomrule
\end{tabular}
\end{table}
\clearpage
